% ============================================================================
% Ridge — KV-cache compression without the attention matrix
% ICLR-workshop skeleton. Four-beat structure from
% notes/paper_submission/paper_story.md.
%
% HOW TO USE THIS IN OVERLEAF
%   1. Overleaf gallery -> search "ICLR" -> "Open as Template" (using ICLR 2025;
%      the style barely changes year to year). That project ships
%      iclr2025_conference.sty + .bst + fancyhdr.sty + math_commands.tex.
%   2. Drop THIS main.tex and references.bib into that project (replace its
%      main.tex). Keep the .sty/.bst files from the template.
%   3. The \usepackage / \bibliographystyle below say iclr2025_conference --
%      confirm that matches the .sty filename in the template's file list.
%      Everything else is style-agnostic.
%
% PAGE LIMIT: default to the tightest (4 pages body; refs/appendix excluded).
% Easier to expand a tight draft than to hack a long one down under deadline.
%
% AI-USE GROUND RULES (paper_story.md): AI is editor/translator, NOT author.
%   - Never let it invent results, numbers, or citations. Every \cite below with
%     a "% VERIFY" tag is a placeholder key you must confirm by hand.
%   - Disclose AI use if the venue asks (usually a one-line statement).
% ============================================================================
\documentclass{article}

\usepackage{iclr2025_conference}   % <- match the template's actual filename
% \input{math_commands.tex}        % ships with the ICLR template; uncomment

\usepackage{booktabs}
\usepackage{graphicx}
\usepackage{amsmath,amssymb}
\usepackage{url}

\title{Ridge: Query-Aware KV-Cache Compression\\ Without the Attention Matrix}

% Keep authorship easy to anonymize (venue may be double-blind). De-anonymize
% only once you confirm the workshop is single-blind.
\author{Anonymous\thanks{Author list withheld for review.}}

\begin{document}
\maketitle

\begin{abstract}
% ONE paragraph, written LAST. The one-sentence claim, then the two-axis result.
% Draft (replace with final numbers):
% Ridge is a KV-cache compressor that selects tokens from the geometry of the
% keys plus a query-importance signal, so it stays competitive with
% attention-based compressors WITHOUT ever materializing the attention matrix --
% making it compatible with FlashAttention. We show Ridge matches the strongest
% attention-based baseline (Compactor) on accuracy at matched compression ratios
% on RULER-16k, while remaining FlashAttention-compatible, which Compactor is not.
% An ablation isolates the contribution of each component (leverage, omega,
% value-weighting, gamma).
\textit{TODO: write last, once Result 1 + ablation numbers are final.}
\end{abstract}

% ===========================================================================
% BEAT 1 — PROBLEM + IDEA  (the depth of the paper)
% ===========================================================================
\section{Introduction}
\label{sec:intro}
% - KV cache is the long-context memory/bandwidth bottleneck; compression = "which
%   tokens do we drop?"
% - Idea, two intuitions: (1) keep geometrically-unique keys (ridge leverage);
%   (2) unique != useful, so add a query-importance signal omega; weight by value
%   substance.
% - Headline clause up front: query-AWARE but never materializes attention =>
%   composes with FlashAttention.
% - Contributions bullet list (mirror the 4 results).

\section{Why leverage alone is not enough}
\label{sec:motivation}
% The intellectual core. Ridge-leverage-unique keys != keys the model attends to.
% Forward-reference the insight figure (Result 4). This motivates omega + gamma.
% See Figure~\ref{fig:leverage-vs-attn}.

\section{Method}
\label{sec:method}
% Compressed version of notes/ridge_explained.md. Define: ridge leverage score,
% omega (query-importance), fixed-envelope combination, value weighting.
% State matched-budget semantics: keep int(T*(1-r)) tokens per head.

% ===========================================================================
% BEAT 2 — POSITIONING + BASELINES
% ===========================================================================
\section{Related Work and Positioning}
\label{sec:related}
% Closest neighbor = Compactor \cite{compactor}: also mixes a leverage-style
% signal with attention, BUT uses materialized attention scores -> incompatible
% with FlashAttention \cite{flashattention}. That gap is our contribution.
% Span the space: SnapKV \cite{snapkv}, knorm, CUR/leverage.

% ===========================================================================
% BEAT 3 — RESULTS  (figures/tables FIRST, prose narrates them)
% ===========================================================================
\section{Experiments}
\label{sec:experiments}
% Setup: Llama-3.1-8B-Instruct, RULER-16k (main), ratios kept 75/50/25% =
% pruned 0.25/0.5/0.75, accuracy at matched kept-token budget per head.

\subsection{Result 1 --- Accuracy vs.\ baselines (main claim)}
\label{sec:res-accuracy}
% ALWAYS include Compactor. Honest tie + practical win beats a suspicious lead.
\begin{figure}[t]
  \centering
  % \includegraphics[width=0.7\linewidth]{figs/accuracy_vs_ratio.pdf}
  \fbox{\parbox[c][3.5cm][c]{0.7\linewidth}{\centering TODO: accuracy-vs-ratio
    line plot. Ridge vs Compactor / SnapKV / knorm / CUR. \textbf{Not yet run}
    --- baselines are commented out in \texttt{sweep.yaml}.}}
  \caption{Accuracy vs.\ compression ratio on RULER-16k, Llama-3.1-8B.}
  \label{fig:accuracy}
\end{figure}

\subsection{Result 2 --- Cost / FlashAttention compatibility}
\label{sec:res-cost}
% Because we tie on accuracy, this is core, not nice-to-have.
\begin{table}[t]
  \centering
  \caption{Compatibility and cost. \textbf{Kernel/H100 numbers pending.}}
  \label{tab:cost}
  \begin{tabular}{lccc}
    \toprule
    Method & Needs attn.\ matrix? & FlashAttn-compatible? & Latency @ long ctx \\
    \midrule
    Compactor & Yes & No  & TODO \\
    SnapKV    & Yes & No  & TODO \\
    Ridge (ours) & \textbf{No} & \textbf{Yes} & TODO \\
    \bottomrule
  \end{tabular}
\end{table}

\subsection{Result 3 --- Ablation (backbone)}
\label{sec:res-ablation}
% Doubles as the method explanation. gamma sweep DONE on LongBench; the
% leverage-only / +omega / +value-weight rungs still to run.
\begin{table}[t]
  \centering
  \caption{Component ablation. Each row adds one piece. \textit{gamma sweep done;
    other rungs pending.}}
  \label{tab:ablation}
  \begin{tabular}{lc}
    \toprule
    Configuration & Accuracy \\
    \midrule
    Leverage only            & TODO \\
    \quad + $\omega$ (query importance) & TODO \\
    \quad + value weighting  & TODO \\
    \quad + $\gamma$ (full Ridge) & TODO \\
    \bottomrule
  \end{tabular}
\end{table}

\subsection{Result 4 --- Unique $\neq$ attended (insight)}
\label{sec:res-insight}
\begin{figure}[t]
  \centering
  \fbox{\parbox[c][3.5cm][c]{0.6\linewidth}{\centering TODO: scatter /
    correlation of ridge-leverage score vs.\ true attention received, per
    key/layer. Coverage logging is wired (\texttt{measure\_coverage}); figure
    not yet produced.}}
  \caption{Ridge-leverage score vs.\ true attention received.}
  \label{fig:leverage-vs-attn}
\end{figure}

% ===========================================================================
% BEAT 4 — FUTURE WORK
% ===========================================================================
\section{Discussion and Future Work}
\label{sec:future}
% Triton kernels for the systems win (score calc + Gram matmuls on GPU).
% Hybrid-model generality (NemotronH) as a short bonus / appendix.

\section*{Reproducibility \& AI-use statement}
% One line if the venue asks. Every claim/number/citation verified by the authors.

\bibliography{references}
\bibliographystyle{iclr2025_conference}   % <- match the template's .bst

\appendix
\section{Additional benchmarks: LongBench}
\label{app:longbench}
% Ridge gamma-ablation + Full baseline on LongBench (16 English tasks) is
% already DONE (24 cells). Lands here as the "does it generalize" evidence.

\section{Hybrid models}
\label{app:hybrid}
% NemotronH: only a few layers carry a KV cache; Ridge still works. Short.

\end{document}
